\subsection{Checkpoint comparisons need matched settings}\label{sec:checkpoint}
% STRENGTHEN CHECKPOINT: Replace the confounded comparison only after the complete
% official-code checkpoint x guidance matrix has been verified.



Our first campaign used \fluxdev, on which \dype{} and \constancy{} produced a dense halftone lattice over the $4096^2$ frame and \ntk{} a featureless blur. An official-code \dype{} check reproduced the lattice on dev prompt 0, seed 42, at guidance 3.5; the corresponding \krea{} run at guidance 4.5 had no visible lattice (\cref{fig:checkpoint}). This rules out our port as a necessary cause of that particular failure, but the guidance difference prevents isolating the checkpoint. \dype's main reported experiments use \krea{}~\citep{issachar2025dype}, as do the official scripts we inspected. We used \krea{} for all held-out comparisons and retain the \fluxdev{} pilot as descriptive evidence only.

\subsection{\sega: paper versus code}\label{sec:papercode}

With the printed law (\cref{eq:sega,tab:papercode}), the allocation is unbounded, and at $4096^2$ the physical bin map sends every pair with a wavelength above \N{c@tokside} tokens, \N{c@collapsedpairs} of the \N{c@pairsperaxis} per axis, to the same frequency bin. The released law scales the spread by 0.15 and clamps the result at 1, and its logit gain stays between $\NmechSegaGainMin$ and $\NmechSegaGainMax$ on our held-out runs. In the eight-prompt dev pilot (descriptive only; \cref{app:pilot}), our implementation of the printed equations scores worse than the released code on every metric; freezing the printed allocation into a static tilt destroys the images, while freezing the released allocation costs much less (\cref{fig:papercode}). Readers who reimplement \sega{} from the paper will not reproduce the released method, and comparisons against \sega{} should state which version they ran.

\subsection{Pitfalls in 4K evaluation}\label{sec:pitfalls}\label{sec:spectra}

\paragraph{224-pixel scorers miss the finest octaves}
CLIP, PickScore and ImageReward see a 224-pixel view, whose Nyquist frequency is \N{nyq@cpp} cycles per original pixel ($2^{\N{nyq@log}}$); bands \N{nyq@firstband} to 10 of our 11 octave bands lie entirely above it. The methods differ in these bands (\cref{fig:spectrum_main,tab:octave}), which may help explain the reward-model ties alongside a full-frame MUSIQ gap. Resizing can retain indirect effects or alias fine patterns, but these scorers cannot directly resolve the original fine detail and are insufficient alone for 4K quality comparisons. \Cref{app:full} gives the bandwise values.

\paragraph{Attention kernel}
Attention backend and GPU can dominate method-level timing differences (\cref{app:compute}); comparisons must fix both. Our records did not log the dispatched kernel or a hash of the loaded weights. We reconstructed the configured attention policy and transformer identity from code revisions and task scripts; neither is a direct runtime record.
% STRENGTHEN QWEN: Insert complete within-family Qwen results here or in a dedicated
% appendix, including every registered arm and metric; do not compare raw rewards
% between families as checkpoint effects.
